\section{引言}
\label{sec:intro}

自动化系统已经在决定哪些帖子被删除、哪些账户被封禁、哪些交易被冻结；越来越多的治理提案还打算将此类决定交由在极少或没有人类参与的情况下运行的 AI 系统作出~\cite{gorwa2020algorithmic,gillespie2018custodians,citron2008technological}。
其实际障碍一直在于成本与延迟：对每条消息进行推理的生成式语言模型过慢、过贵，无法部署在每一次输入与输出之前。
2026 年 9 月 15 日，TypeSafe AI 发布了 \jev{}，该模型完全不生成文本，从而消除了这一障碍~\cite{almeida2026introducing,typesafe_intro}。
给定一个输入和调用方声明的问题，它返回在所列选项中的选择、按评分细则给出的分数，或某一陈述为真的概率，并各自附带厂商声称经过校准的概率~\cite{typesafe_primer,typesafe_models}。
由于输出是一个数值而非文字，普通代码即可对其设定阈值并据此行动，其延迟和价格仅为一次生成式调用的一小部分。
我们将此类系统称为\emph{类型化决策模型}；厂商自己的术语“System One”模型借用了快速、直觉式判断的双过程隐喻~\cite{kahneman2011thinking}。
数日之内，具有相同接口的开放权重模型相继出现；截至 2026 年 10 月 1 日，已有数十项研究对其进行了评估。

本文围绕 PoL 治理，即依据爱2证明（\pol{}）规范进行的治理，探讨类型化决策模型能否以及在何种条件下充当其自动化保障机制；由于该规范使机器判断异常具体，其答案也关乎更广泛的 AI 治理。
\pol{} 是一份开放的、带版本管理的治理文本~\cite{pol2text,pol2en}，它在现有大语言模型（LLM）内部设置\emph{内置的共识治理层}，其中的\emph{安全阀}对输入和输出进行审计与拦截，使恨语“在必要情况下”被阻断（\clause{5.4.1}）。
\Cref{sec:paradigms} 通过与监管、平台、训练期、护栏和 DAO 等治理范式的比较，说明 \pol{} 的哪些特性使这一问题变得具体，以及每一特性的代价。中国 2023 年的《生成式人工智能服务管理暂行办法》~\cite{cac2023genai} 同样将拦截义务、其所适用的标准以及附着于决定的义务结合在一起；\pol{} 之所以值得研究，并非因为其独特，而是因为其构念与拦截点（每一次输入与输出）均在可引用的条款中写明，可从中推导出可检验的要求，而安全阀如何运作则未作规定，同时也因为其缺口清晰可见。
其人类公共福利的治理协议规定，由 AI 组织 NaturalDAO 负责无需人类投票的福利决策、异常检测、可验证的决策链路以及争议裁决，而人类监督“不直接干预”（\art{2}--\art{10}）。
若将这些条款视为工程要求，它们所要求的恰恰是类型化模型所能提供的快速、结构化、可记录的判断；同时，它们也清楚地揭示了其社会利害：这些判断将限制言论、分配福利并解决争议。
因此，本文的研究问题是：
\begin{quote}
\emph{类型化决策模型能在多大程度上、在何种条件下承担诸如 \pol{} 安全阀这类受条款约束的 AI 保障机制所要求的判断？这些证据对此类规范应如何撰写与实现意味着什么？}
\end{quote}

\paragraph{检测器还是裁决者。}
本文的结论取决于两种角色的区分。
\emph{检测器}产生一个有记录的信号，用于排序、确定优先级或触发审查，或仅触发有时限且完全可逆的行动。
\emph{裁决者}是指任何这样的自动化组件或组合：其输出在没有同步人工确认的情况下，决定一项在短暂有限期限之外限制某人言论、资源或账户的结果，或决定一项在申诉时无法被完全撤销的结果。
自动\emph{放行}不限制任何人，因而仍属于检测器角色，但它对受被放行内容伤害的任何人都有实质后果；因此我们将其视为经过审慎考虑、受审计的例外：其漏检率与误阻断率一样是受治理的量，放行决定通过随机抽样进行审计。
本文仅在此意义上使用\emph{裁决者}一词；凡文献中称“LLM judges”或“LLM-as-a-judge”之处，除论文标题外，我们均写作 \emph{LLM 评估器}。

\paragraph{论点。}
本文的论点是有条件的、比较性的。
截至 2026 年 10 月 4 日检索到的证据支持在附加条件下将类型化决策模型用作检测器：它们在英文基准上对风险内容的排序效果良好，成本仅为 LLM 评估器的一小部分；在一项基于模拟数据的预注册研究中，仅当一条独立构建的规则同时认可时才接受其低风险判定，这一做法相对于前沿 LLM 未造成准确率损失，但伪装攻击仍得以通过。
证据不支持将其用作裁决者：阈值无法在不同情境之间迁移，判定可能取决于选项的名称而非其定义，插入的观点可以操纵它们，明确的“不知道”选项在其为正确答案时很少被选择，而“以上皆非”的选择不一致；并且在评分细则打分任务中，低成本的同类评估器几乎重复了某一类型化模型置信度最高的全部错误。
在同一要求上对类型化与生成式评估器进行比较时，类型化模型总体上并不更差；但多数失效从未经过比较，因此这些证据既不能表明失效为类型化模型所特有，也不能表明这些失效为二者所共有。
尚无研究测量 \pol{} 所需的构念，即对中文中的爱语、恨语与不在场状态进行三值判断，因此即便是关于检测器的论断也属于外推，且每一项发现的确定性均为低或极低（\cref{tab:keyfindings}）。
因此，对于 \pol{} 本身，本文给出明确结论（\cref{sec:verdict}）：依现有证据，托管的 \jev{} 不适合承担该规范中任何作出决定的功能，即争议裁决、自动伦理验证、逐人量化（依据操纵方面的证据，并且由于尚无研究按群体测量误差，亦基于审慎原则）或对某人言论的最终阻断；对于检测类功能，即安全阀、输出截断以及针对文本的异常检测（若无经过训练的模型，则不适用于数值或结构化信号），它只能作为\cref{sec:design}设计中的留有记录的检测器使用，而 \clause{5.6} 的公共决策评估只能作为一种经审计的汇总测量，并将其测得误差带入每一项结论；这一用法是从英文基准外推而来的。
截至 10 月 8 日的更新检索单独报告于\cref{sec:update}，它未推翻上述任何一项回答：它限定了一项表述（关于分解式决策链路的准确率），将一项发现（预先设定的阈值在新数据上达不到其错误目标）的确定性评级由极低上调为低，并新增了一项研究，使另一发现（升级至更强的推理模型可在不损失准确率的情况下节省成本）有争议。

\paragraph{贡献。}
\begin{enumerate}
  \item \textbf{对一份治理规范的要求分析。} 我们将 \pol{} 中要求或约束机器判断的条款转化为七项可检验的要求 G1--G7（\cref{sec:method}），并以中英双语引用这些条款及本文所引的相关段落（来自 \nPolClauses{} 个条款的段落，见\cref{app:concordance}）；并将该规范置于另外五种治理范式之中，以说明其哪些特性使其条款可被解读为推理期保障机制的可检验要求，以及每一特性的代价（\cref{sec:paradigms}）。
  \item \textbf{结构化的证据评估。} 我们依据这些要求评估 \nPapers{} 项研究和 \nGrey{} 项灰色文献：凡与某项要求相关的文献，均由 AI 智能体在作者指导下全文阅读、质量评价，并采用对称的编码手册进行编码；随机抽取的样本由另一 AI 编码者进行盲法双重编码；关键结论均附有确定性评级（\cref{sec:method,sec:evidence}）。
  \item \textbf{条款层面的政策分析。} 这些证据凸显了该规范内部的五项张力，涉及情绪识别、二值判定、封闭模型下的透明性、AI 裁决中的非干预原则以及逐人评分；我们将其与欧盟《人工智能法》、《通用数据保护条例》（GDPR）、《数字服务法》及内容审核规范进行比较（\cref{sec:synthesis}）。
  \item \textbf{附带验收标准的参考设计。} 我们提出一种保障机制，其中类型化模型作为有记录、三值、经本地校准的检测器，而具有实质后果的决定最终由人作出；同时给出任何模型在部署前都应满足的暂定通过标准（\cref{sec:design}）；并逐项功能给出明确结论，说明依现有证据，托管的 \jev{} 是否适合作为 \pol{} 的工具（\cref{sec:verdict}）。
\end{enumerate}
证据评估是回答政策问题的方法：我们并不尝试对类型化决策模型进行全面综述，仅在研究与 G1--G7 相关的范围内加以讨论。

\paragraph{与已有工作的关系。}
与本文最接近的已有工作是对最早 28 篇类型化模型论文（9 月 19 日至 24 日）的审查，该审查发现，相较于从生成式模型中读取标签概率，类型化读出尚未显示出独立的准确率优势，并提出了一份包含 14 项内容的评估核查清单~\cite{arx32160}。
本文覆盖时间更晚、范围更大的窗口，包括关于观点攻击、弃权、概率一致性和门控级联的首批研究；我们全文阅读了每一项文献；并且我们围绕一份治理规范及其监管背景而非围绕模型评估来组织证据。
本文的评估核查清单在~\cite{arx32160} 的基础上增加了通过标准（\cref{tab:checklist}）。
此前发布的一份数据集~\cite{survey01} 主要依据摘要对研究进行了总结；本文取代了该版本，\cref{app:corrections} 列出了实质性更正。

\paragraph{范围。}
我们以 \pol{} 条款作为起始规范，探讨能否以及如何依据现有证据实现这些条款。
凡证据涉及某一条款能否按原文实现之处，\cref{sec:synthesis} 指出相应张力并提供若干选项，其中至少一项保留原文；我们不评价 \pol{} 更广泛的规范性或政治性承诺。
作者是 \pol{} 的贡献者之一（见\hyperref[sec:statements]{立场声明与利益冲突}）。
本文结论适用于 \texttt{jev-1.13}（或部分研究中所使用的未报告版本的托管模型）以及 2026 年 10 月 1 日可获得的开源模型。
